\documentclass[11pt,aspectratio=169]{beamer}

\usetheme{Madrid}
\usecolortheme{seahorse}

\usepackage[utf8]{inputenc}
\usepackage{amsmath,amssymb}
\usepackage{xcolor}
\usepackage{booktabs}
\usepackage{tikz}
\usetikzlibrary{arrows.meta,positioning,fit,backgrounds}

\setbeamertemplate{navigation symbols}{}
\setbeamertemplate{footline}[frame number]

\title{Beyond the Syllabus}
\subtitle{Sharper Hardness, Solvers That Work Anyway, and Why Lower Bounds Are
Hard to Prove\\
\small Week 12 Supplement -- TOAE209/TOAH209: Design and Analysis of Algorithms}
\author{Foreign Trade University}
\date{}

\begin{document}

% ---------------------------------------------------------------
\begin{frame}
  \titlepage
\end{frame}
% ---------------------------------------------------------------

\begin{frame}{Outline}
  \tableofcontents
\end{frame}

% =================================================================
\section{Motivation}
% =================================================================

\begin{frame}{Why look beyond the syllabus?}
  \small
  \begin{itemize}
    \item This week you learned that \textsc{3-SAT}, \textsc{Vertex Cover},
    \textsc{Clique}, \textsc{Hamiltonian Cycle}, \textsc{Subset Sum}, and
    \textsc{3-Coloring} are all NP-complete -- and a poly.\ algorithm for one would give
    one for \emph{all} of them.
    \item ``NP-complete'' is not the end of the story -- it raises three sharper
    questions, all active research right now:
    \begin{enumerate}
      \item \textbf{How exponential, exactly?} $2^{0.001n}$ and $2^{n}$ are both
      ``non-polynomial,'' but wildly different in practice.
      \item \textbf{If it's exponential in theory, why do real SAT solvers routinely
      solve formulas with millions of variables?}
      \item \textbf{Why is it so hard to even \emph{prove} these lower bounds?} We have
      overwhelming \emph{evidence} P $\neq$ NP, but no proof -- and a precise reason our
      current techniques cannot supply one.
    \end{enumerate}
  \end{itemize}
  \begin{alertblock}{The throughline}
    Reductions do more than classify problems as ``hard'' -- they let a hardness
    \emph{assumption} about one problem (like 3-SAT) transfer precise, quantitative
    consequences to every problem it reduces to.
  \end{alertblock}
\end{frame}

% =================================================================
\section{How Exponential, Exactly? The Exponential Time Hypothesis}
% =================================================================

\begin{frame}{Recap: what NP-completeness does \emph{not} tell you}
  \begin{itemize}
    \item Cook--Levin plus this week's reductions prove: \emph{if} P $=$ NP, \emph{then}
    3-SAT (and Vertex Cover, Clique, \dots) has a polynomial algorithm.
    \item They say nothing about \emph{how exponential} the best algorithm is if
    P $\neq$ NP. Brute-force 3-SAT is $O(2^n)$; is that close to necessary, or could a
    clever algorithm reach $O(1.01^n)$, or even $O(2^{\sqrt{n}})$?
    \item This is exactly the kind of question this course trains you to ask about
    \emph{every} algorithm's running time -- we are just asking it about the
    \emph{lower} bound side now.
  \end{itemize}
\end{frame}

\begin{frame}{The Exponential Time Hypothesis (Impagliazzo \& Paturi)}
  \begin{block}{Impagliazzo \& Paturi, ``On the Complexity of $k$-SAT'' (CCC 1999;
  \emph{J.\ Comput.\ Syst.\ Sci.}, 2001)}
    For $k \ge 3$, let $s_k = \inf\{\delta : k\text{-SAT solvable in } O(2^{\delta n})
    \text{ time}\}$. The \textbf{Exponential Time Hypothesis (ETH)} conjectures
    $s_3 > 0$: 3-SAT has \emph{no} $2^{o(n)}$ algorithm.
  \end{block}
  \begin{itemize}
    \item A companion \textbf{Sparsification Lemma} (Impagliazzo, Paturi \& Zane, 2001)
    shows any $k$-SAT instance can be re-expressed as $2^{\varepsilon n}$ instances with
    only $O(n)$ clauses each -- the technical engine that makes ETH a \emph{robust},
    reduction-friendly assumption rather than a fact about one particular encoding.
    \item ETH is \emph{stronger} than P $\neq$ NP (it rules out sub-exponential time,
    not just polynomial time) and is now a standard assumption for proving
    \textbf{quantitative} lower bounds via this week's reduction technique.
  \end{itemize}
\end{frame}

\begin{frame}{Why this matters for problems from this week's lecture}
  \begin{itemize}
    \item ETH-based reductions give lower bounds like: \textsc{Vertex Cover} on an
    $n$-vertex graph has no $2^{o(n)}$ algorithm, and \textsc{Independent Set} /
    \textsc{Clique} inherit the same bound through the reductions you proved this week
    (Lemma: Independent Set $\equiv$ Vertex Cover; Independent Set $\le_p$ Clique).
    \item For \textsc{Subset Sum}: the pseudo-polynomial $O(nW)$ DP you may see next
    week is \emph{not} an accident of a weak algorithm -- fine-grained results (building
    on this same ETH/SETH line) show truly-polynomial-in-$n$-and-$\log W$ algorithms
    would violate standard hardness assumptions.
    \item \textbf{The upshot:} once you have proved $X$ is NP-complete this week, ETH
    lets you say something much sharper than ``no poly algorithm'' -- it predicts
    \emph{how much} exponential blow-up is unavoidable, which is exactly the number you
    need to decide whether brute force, or the FPT/branch-and-bound tools of Week 13,
    is the realistic plan.
  \end{itemize}
\end{frame}

% =================================================================
\section{NP-Complete in Theory, Solved in Practice: Modern SAT Solvers}
% =================================================================

\begin{frame}{The puzzle: exponential worst case, routine industrial use}
  \begin{itemize}
    \item ETH says 3-SAT has no sub-exponential \emph{worst-case} algorithm. And yet
    modern SAT solvers are used daily in hardware verification, software model
    checking, and scheduling, on formulas with \textbf{millions} of variables.
    \item The resolution: worst-case exponential and \emph{practically fast on the
    instances that actually arise} are compatible. Conflict-Driven Clause Learning
    (CDCL) solvers combine backtracking search (this week's topic) with learned
    clauses that prune huge parts of the search space the moment a conflict repeats a
    pattern.
    \item Progress on this engineering problem is measured every year at the
    \textbf{SAT Competition} -- a genuine, ongoing benchmark of how far practice has
    outrun the worst-case theory.
  \end{itemize}
\end{frame}

\begin{frame}{SAT Competition 2024--2025: who is winning, and how}
  \small
  \begin{center}
  \begin{tabular}{@{}l l p{5.4cm}@{}}
    \toprule
    Year & Main sequential track winner & Note \\
    \midrule
    2024 & \textbf{Kissat} & a from-scratch C port/redesign of CaDiCaL \\
    2025 & \textbf{CaDiCaL-SC2025} & PAR-2 $2327.00$, $161$ solved; Kissat-VSA 2nd
    ($2335.54$, $160$ solved) \\
    \bottomrule
  \end{tabular}
  \end{center}
  \begin{itemize}
    \item Source: SAT Competition 2025 official results
    (\texttt{satcompetition.github.io/2025}) and the SAT Competition 2024 medal table
    (\texttt{cca.informatik.uni-freiburg.de}).
    \item Striking detail: \emph{all three} medalists in the 2025 main sequential track
    were Kissat-based variants (augmented with extra heuristics) -- the CDCL
    architecture from this week's backtracking discussion, refined for 25+ years, is
    still the state of the art.
    \item This is the same lesson as the exercises you just did: brute-force
    backtracking plus the \emph{right} pruning rule is the difference between
    ``theoretically exponential'' and ``solves the instance in milliseconds.''
  \end{itemize}
\end{frame}

% =================================================================
\section{Why Is Proving These Lower Bounds So Hard?}
% =================================================================

\begin{frame}{We believe P $\neq$ NP. Why can't we prove it?}
  \begin{itemize}
    \item Every standard technique for proving a Boolean function is \emph{hard to
    compute with small circuits} has, historically, worked by exhibiting a
    \textbf{property} of hard functions that is: (i) \emph{constructive} (checkable in
    time polynomial in the function's truth table), (ii) \emph{large} (true for most
    functions), and (iii) \emph{useful} (false for the easy functions you're separating
    from).
    \item \textbf{Razborov \& Rudich (1994/1997), ``Natural Proofs.''} Any proof
    technique with all three properties would, as a side effect, break standard
    pseudorandom generators / one-way functions -- objects that cryptography needs to
    exist. So either cryptography is broken, or \emph{no ``natural'' proof can ever
    separate P from NP.} This is the \textbf{natural proofs barrier}.
  \end{itemize}
\end{frame}

\begin{frame}{2026: the same barrier, for data structures}
  \begin{block}{Loff, Kouck\'{y}, Molli \& Saks, ``The Natural Proofs Barrier against
  Data-Structure Lower-Bounds'' (STOC 2026; DOI: 10.1145/3798129.3800843)}
    The natural-proofs obstruction is not special to circuit complexity: the same
    style of argument shows that ``natural'' techniques are also insufficient to prove
    strong \textbf{data-structure} lower bounds (e.g.\ for dynamic connectivity-style
    problems closely related to this week's reductions), under analogous
    cryptographic-hardness assumptions.
  \end{block}
  \begin{itemize}
    \item Why this belongs in a week on NP-completeness: it is a direct descendant of
    the verification-vs-search asymmetry this week is built on -- it says the
    \emph{proof techniques themselves} face a search-vs-verify-style obstacle,
    one level up.
    \item Practical upshot for you: when an exercise asks you to prove a lower bound,
    ``try harder with the natural strategy'' sometimes has a hard mathematical ceiling
    -- this is an active 2020s research frontier, not a solved problem.
  \end{itemize}
\end{frame}

% =================================================================
\section{Summary}
% =================================================================

\begin{frame}{Summary}
  \footnotesize
  \begin{enumerate}
    \item \textbf{Sharper hardness.} The Exponential Time Hypothesis (Impagliazzo \&
    Paturi, 1999/2001) upgrades ``no polynomial algorithm'' to a quantitative
    $2^{\Omega(n)}$ prediction, transferred to Vertex Cover/Clique/Subset Sum by exactly
    the reductions you proved this week.
    \item \textbf{Theory vs.\ practice.} Worst-case exponential and
    ``routinely solves million-variable industrial instances'' coexist: SAT
    Competition 2024--2025 medalists (Kissat, CaDiCaL) are still backtracking search
    plus learned-clause pruning -- this week's idea, engineered for 25 years.
    \item \textbf{Why lower-bound proofs are themselves hard.} The natural proofs
    barrier (Razborov--Rudich, 1994/1997) blocks the standard circuit-lower-bound
    recipe unless cryptography is broken; a STOC 2026 result (Loff, Kouck\'{y}, Molli,
    Saks) shows the same obstruction reaches data-structure lower bounds.
  \end{enumerate}
  \begin{alertblock}{Looking ahead}
    Week 13 (Kleinberg \& Tardos, Ch.\ 10) takes NP-completeness as a
    \textbf{starting point} rather than a dead end: fixed-parameter tractable
    algorithms confine the exponential blow-up to a small parameter $k$, PSPACE
    captures games and quantified logic, and branch-and-bound turns today's
    backtracking into a principled search strategy with provable pruning.
  \end{alertblock}
\end{frame}

\end{document}
